\begin{frame}{왜 ``보상''만으로는 부족할까: 지연 강화}
  \textbf{지도학습 vs 강화학습: 보상 시점 차이}
  
  \vskip0.6em
  \textbf{지연 강화: RL의 두 번째 핵심 축}
  \begin{block}{시간차(TD) 학습}
    ``지금의 추정치를, \textbf{한 스텝 미래의 추정치 + 즉시 보상}으로
    보정한다'' --- 1988년 서턴이 정식화. $\Rightarrow$ \textbf{6주차}의 주제.
  \end{block}

\note{\begin{itemize}
\item 지도학습의 라벨은 \textit{그 사진 바로 옆}에 붙어 있음.
  강화학습에서는 보상이 \textbf{늦게} 옴.
\item 바둑에서 \textbf{50수 전}에 둔 수가 50수 뒤의 승패로
      드러남.
\item 로봇이 걷다가 \textbf{100스텝 뒤}에 넘어졌을 때,
      ``어느 스텝의 행동이 원인인가''를 \textbf{에이전트가 스스로}
      추론해야 함.
\item 이 \textbf{신호가 시간적으로 벌어져 있다}는 사실이 RL을 구분하는
  \textit{두 번째} 핵심 축. 이를 푸는 도구:
\end{itemize}}
\end{frame}
